%%%%%
%
\section{Transformation d'un temps GPS en temps UTC (version automatique asynchrone)}


\itemstyle{Op\'erateur~:} \operfmt{GPS\_to\_UTC\_irreg}
\index{GPS\_to\_UTC\_irreg}
\label{OP-GPS_to_UTC_irreg}

%%%%%OBJET%%%%
\itemstyle{Objet~:}
Conversion en UTC de temps exprim�s dans le syst�me GPS (en secondes
depuis le d�but de la semaine, d�cal� d'une valeur constante avec le
TAI et d'une valeur variable avec l'UTC).

Le r�sultat est exprim� en secondes depuis 0h00, selon la
convention g�n�rale AIDA.

Cet op�rateur fonctionne automatiquement en utilisant une
valeur du d�calage GPS-UTC fournie par
le r�cepteur (sur une base de temps qui n'est pas forc�ment synchrone
avec les mesures de position).


%%%%%ARGUMENTS%%%%
\bigskip
%%%%%ENTREES%%%%
\begin{tabarg}
\multicolumn{3}{l}{\titretab {Entr\'ees}}\\
\hline
$T_{GPS}$ & Vecteur, \texttt{float} & Temps exprim� dans le syst�me GPS (second)\\
$Decal$ & Vecteur, \texttt{float} & Information sur le d�calage GPS/UTC tel
que fournie par le syst�me Trimble/Bancom utilis� � SAFIRE (second)\\
$N_{Decal}$ & Coeff., \texttt{U\_Int} & Nombre d'�chantilllons de $Decal$ \\
\end{tabarg}

\medskip
%%%%%SORTIES%%%%
\begin{tabarg}
\multicolumn{3}{l}{\titretab {Sorties}}\\
\hline
$T_{UTC}$ & Vecteur, \texttt{double} & $T_{GPS}$ exprim� dans le syst�me UTC (second)\\
\end{tabarg}


%%%%%FORMULE%%%%
\itemstyle{Formule}

%% GPS_to_UTC_irreg
  Recherche des valeurs non-nulles de $Decal$, et
v�rification qu'elles sont identiques. Si c'est le
cas, on d�finit $D$ comme cette valeur.

  Sur le premier $T_{GPS}$ non-nul, calcul d'un offset : 
\[
offset = 86400 . [(T_{GPS}-D) / 86400]
\]
Puis :
\[
   T_{UTC} = T_{GPS} - D - offset
\]




\itemstyle{Commentaire~:}
Si un vecteur des diff�rences GPS-UTC n'est pas disponible ou
inutilisable, il faut utiliser l'op�rateur
\hyperref[OP-GPS_to_UTC]{\operfmt{GPS\_to\_UTC}} avec un argument explicite.

%%%%%AUTEUR%%%%
\itemstyle{Auteur~:}
CNRM/GMEI/TRAMM

%%%%%REFERENCE%%%%
\itemstyle{R\'ef\'erences~:}
\textit{GPS, theory and practice}, par B. Hofmann-Wellenhof, H. Lichtenegger,
and J. Collins \cite{Hofmann-GPS}
